\section{The upper constant 12 of the RMS law}
\label{app:upper-12}

The proof of \cref{thm:main} in \cref{sec:minimax-proof} gives the upper
constant $\sqrt{8\cdot12666}<319$, and its formal version gives $512$.
This appendix proves the constant $12$ for one distribution-blind family.
The family runs the calibrated batch of one level, as in
\cref{sec:minimax-proof}, with every call that would use no observation
padded by one, as in its formal version. Its calibration denominator is
$B/4$, which need not be a power of two, so the batch uses $B$ minus the mean
cost of one call in expectation, in place of as little as about $B/2$ minus
that cost. The other
improvements are sharper bounds on the second moment and the cost of one call
of the cutoff estimator.

\begin{theorem}\label{thm:upper-12}
For each $B\ge1$ there is one algorithm $\cA_B$, depending only on $B$,
such that for every probability law $\mu$ on $(0,1/2]$ and every sign
function $q$,
\begin{equation}\label{eq:upper-12}
 \E_{\mu,q}T\le B,\qquad
 \bigl(\E_{\mu,q}(\widehat\theta-\theta)^2\bigr)^{1/2}\le12\,\om(B^{-1/2}).
\end{equation}
In particular $R_\mu(B)\le12\,\om(B^{-1/2})$. With \cref{thm:sharp-rms},
the two sides of \eqref{eq:main-rms} are within the factor $72$.
\end{theorem}

The constant $12$ is not claimed optimal; \cref{sec:upper-12-scope}
describes the remaining slack.

\paragraph{Formal scope.}
\upperleanref{UpperScale/Full.lean}{blind\_upper\_12} proves
\eqref{eq:upper-12} for one budget-indexed family chosen before the gap
law, against the same experiment, budget, risk and modulus definitions as
the formal statement of \cref{thm:main}, and
\upperleanref{UpperScale/Full.lean}{minimax\_upper\_12} proves the
minimax bound. The Lean proof follows the route below, and the Lean kernel
checks the numerical certificates that the text summarizes: the offset
bounds of \cref{lem:upper-12-offset-cost,lem:upper-12-offset-moment} and
the drift steps of \cref{lem:upper-12-drift}.
\Cref{sec:upper-12-scope} lists the modules and the audit.

\subsection{The algorithm and the reduction to one call}

Recall the cutoff estimator of \cref{sec:classifier}. For a dyadic cap
$M$, the classifier with cap $M$ checks the prefixes of lengths
$n=1,2,4,\dots$ below $M$ against the thresholds $t_{n,M}$ of
\eqref{eq:threshold}, and $\tau_M$ is its number of observations. One call
at cutoff $L$, with $N=2^L$, draws the level $J$ with
$\Prob(J=\ell)=p_\ell=(1-r)r^\ell$, $r=2^{-1/2}$, and returns $Y_L$ of
\eqref{eq:Y-cut}. Its padded cost $C_L$ is the actual work
$\mathcal W_L$ when $J\le L$, and one observation of a fresh coin when
$J>L$. Then $1\le C_L\le N$, and
$\E[C_L\mid d,S]=\sum_{\ell\le L}p_\ell\E[W_\ell\mid d,S]+r^{L+1}$.

For $\lambda\ge1$, the \emph{batch at level $L$ and scale $\lambda N$}
applies \cref{lem:calibration} to $(Z,C)=(Y_L,C_L)$ with the calibration
denominator $\lambda N$ in place of the $N$ of that lemma, which is allowed
because $C_L\le N\le\lambda N$. Given its cost $C_L$, each pilot call is
marked with probability $C_L/(\lambda N)$; at the second mark, after $K$ pilot
calls, the batch averages $K-1$ fresh production calls. For $B<4$, $\cA_B$
observes nothing and outputs $0$, as the formal family does; any fixed output
in $[0,1]$ would do. For $B\ge4$, let $L\ge0$ satisfy
$4\cdot2^L\le B<8\cdot2^L$, so that the cap is $N=2^L$, and put
$\lambda=B/(4\cdot2^L)\in[1,2)$. Then $\cA_B$ runs the batch at level $L$
and scale $\lambda2^L=B/4$. The algorithm uses neither $\mu$ nor $q$.

Put $\rho_N(d)=\sqrt N$ when $Nd^2\le1$ and $\rho_N(d)=0$
otherwise. The following bounds on one call replace the constants $85$
and $149$ of \cref{lem:universal}.

\begin{lemma}\label[lemma]{lem:upper-12-call}
For every $L\ge0$, $N=2^L$, both signs and $0<d\le1/2$,
\begin{align}
 \E[C_L\mid d,S]&\le\tfrac{26}{5}\min\{\sqrt N,d^{-1}\}
   -\tfrac{21}{5}\rho_N(d),\label{eq:upper-12-cost}\\
 \E[Y_L^2\mid d,S]&\le\tfrac{13}{2}\min\{\sqrt N,d^{-1}\}
   +\tfrac{129}{25}\rho_N(d).\label{eq:upper-12-moment}
\end{align}
\end{lemma}

The regime term $\rho_N$ is needed because for $Nd^2\le1$ the bounds
used here give the second moment only $(6+4\sqrt2)\sqrt N<11.66\sqrt N$,
which exceeds $\tfrac{13}{2}\sqrt N$. In that regime the padded cost is
at most $\sqrt N$, far below $\tfrac{26}{5}\sqrt N$, and the product of
the two bounds is what matters.

\begin{lemma}\label[lemma]{lem:upper-12-batch}
For every $L\ge0$, $N=2^L$ and $\lambda\ge1$, the batch at level $L$ and
scale $\lambda N$ uses at most $4\lambda N$ expected observations and
$\E(\widehat\theta-\theta)^2\le\bigl(1+\tfrac{169}{5\lambda}\bigr)\om(N^{-1/2})^2$,
for every $\mu$ and $q$.
\end{lemma}

\begin{proof}
The work bound is \eqref{eq:calibration} with $\lambda N$ in place of $N$.
Put $h=N^{-1/2}$, $A=\Am(h)=\E_\mu\min\{\sqrt N,D^{-1}\}$ and
$G=\E_\mu\rho_N(D)\ge0$. Averaging \cref{lem:upper-12-call} over $(D,S)$
bounds the variance $V$ and the mean cost $\bar C$ of one call, the $v$ and
$c$ of \cref{lem:calibration}, by
$V\le\tfrac{13}{2}A+\tfrac{129}{25}G$ and
$0\le\bar C\le\tfrac{26}{5}A-\tfrac{21}{5}G$. Hence
\[
 V\bar C\le\tfrac{13}{2}\cdot\tfrac{26}{5}A^2
  +\bigl(\tfrac{129}{25}\cdot\tfrac{26}{5}-\tfrac{13}{2}\cdot\tfrac{21}{5}\bigr)AG
  -\tfrac{129}{25}\cdot\tfrac{21}{5}G^2
 \le\tfrac{169}{5}A^2,
\]
since $\tfrac{129}{25}\cdot\tfrac{26}{5}=\tfrac{3354}{125}<\tfrac{273}{10}
=\tfrac{13}{2}\cdot\tfrac{21}{5}$. By \eqref{eq:calibration} and the bias
bound of \eqref{eq:one-call}, the mean-square error is at most
$\om(h)^2+V\bar C/(\lambda N)$, and $A^2/N=\om(h)^2$ because $\om(h)=hA$.
So it is at most $(1+\tfrac{169}{5\lambda})\om(h)^2$.
\end{proof}

\begin{proof}[Proof of \cref{thm:upper-12}]
For $B<4$, $\om(B^{-1/2})=1$ because $B^{-1/2}>1/2\ge D$, and the
output $0$ has error at most $1$. Let $B\ge4$ and $N=2^L$. By
\cref{lem:upper-12-batch}, the expected work is at most $4\lambda N=B$.
Since $N\le B$, \eqref{eq:scaling} with $c=\sqrt{B/N}$ gives
$\om(N^{-1/2})^2\le(B/N)\om(B^{-1/2})^2=4\lambda\,\om(B^{-1/2})^2$. So the
mean-square error is at most
\[
 4\lambda\bigl(1+\tfrac{169}{5\lambda}\bigr)\om(B^{-1/2})^2
 =4\bigl(\lambda+\tfrac{169}{5}\bigr)\om(B^{-1/2})^2
 <\tfrac{716}{5}\,\om(B^{-1/2})^2,
\]
because $\lambda<2$, and $\tfrac{716}{5}<12^2$. Taking the infimum over
admissible algorithms gives the minimax bound.
\end{proof}

The batch at the dyadic scale $2^L$ with $4\cdot2^L\le B$ would pay the
factor $8$ of \cref{sec:minimax-proof} in the worst case, at $B$ just
below $8\cdot2^L$. A mixture of the batches of levels $L$ and $L+1$, each
at its dyadic scale, also has expected work at most $B$, but pays the factor
$9/2$, which with the bounds of \cref{lem:upper-12-call} gives
$\sqrt{(9/2)\cdot174/5}\approx12.5$. The calibration at scale $B/4$ pays
the factor $4$; the bias stays that of the dyadic level $L$, and its square
is at most $\lambda<2$ times $\om((B/4)^{-1/2})^2$.

\subsection{The cost of one call}

Since $t_{n,2M}>t_{n,M}$, the classifier with cap $2M$ never stops
before the classifier with cap $M$. The coupled work of level $\ell\ge1$
is therefore $W_\ell=\tau_{2^\ell}$, and $W_0=\tau_1=1$.

The survival identity in \eqref{eq:survival} gives
$\E\tau_M=1+\sum_{n<M}n\Prob(\tau_M>n)$, the sum over dyadic $n$.
Surviving the prefix $n$ requires $|\overline X_n-1/2|\le t_{n,M}$. When
$t_{n,M}\le d$, Hoeffding's inequality at the deviation $d-t_{n,M}$
bounds this by $e^{-2n(d-t_{n,M})^2}$, for either sign. Write $n=2^k$,
$M=2^\ell$ and $j=\ell-k\ge1$, and put
\[
 v_j=1+j\log2,\qquad y_k=2^kd^2 .
\]
Then $nt_{n,M}^2=v_j$, so $t_{n,M}\le d$ exactly when $y_k\ge v_j$, and
then $n(d-t_{n,M})^2=v_j\bigl(\sqrt{y_k/v_j}-1\bigr)^2$. The threshold
depends on $\ell$ and $k$ only through $j$, which the proof of
\cref{lem:classifier} does not use. Hence
\begin{equation}\label{eq:upper-12-survival}
 \E\tau_{2^\ell}\le1+\sum_{k<\ell}2^k\pi_{\ell-k}(y_k),\qquad
 \pi_j(y)=\begin{cases}1,&y<v_j,\\
 e^{-2v_j(\sqrt{y/v_j}-1)^2},&y\ge v_j.\end{cases}
\end{equation}

\begin{lemma}\label[lemma]{lem:upper-12-offset-cost}
For $v>0$ and $s\in[1,2]$ put
\[
 F_v(s)=(1+\sqrt2)\sqrt s
  +\sum_{i\ge0}\sqrt{2^is}\,e^{-2v(\sqrt{2^is}-1)^2}.
\]
Then $F_v(s)\le89/20$ for $v\ge169/100$, $F_v(s)\le17/4$ for
$v\ge119/50$, and $F_v(s)\le399/100$ for $v\ge307/100$.
\end{lemma}

\begin{proof}
$F_v$ decreases in $v$, so it suffices to take $v=169/100$, $v=119/50$
and $v=307/100$. The Lean proofs split $[1,2]$ into $11$, $9$ and $29$
pieces, chosen and checked in exact rational arithmetic. On each piece
they bound the first four terms of the series by the monotonicity of each
factor and an eight-term Taylor lower bound for the exponential, and the
remaining terms by $1/80$.
\end{proof}

\begin{proof}[Proof of \eqref{eq:upper-12-cost}]
\emph{The cap bound.} Since $\tau_{2^\ell}\le2^\ell$,
\[
 \E[C_L\mid d,S]\le(1-r)\sum_{\ell\le L}2^{\ell/2}+r^{L+1}
 =\sqrt N-r+r^{L+1}\le\sqrt N .
\]
If $Nd^2\le1$, the right side of \eqref{eq:upper-12-cost} is
$(\tfrac{26}{5}-\tfrac{21}{5})\sqrt N=\sqrt N$. If $1<Nd^2\le16$, then
$\sqrt N\le4/d<\tfrac{26}{5}\min\{\sqrt N,d^{-1}\}$.

\emph{Reordering by $j$.} Let $Nd^2>16$; since $d\le1/2$, then $L\ge7$.
By \eqref{eq:upper-12-survival}, and since
$p_\ell2^k=(1-r)r^j\sqrt{y_k}/d$ for $\ell=k+j$,
\[
 d\,\E[C_L\mid d,S]\le d\Bigl(\sum_{\ell\le L}p_\ell+r^{L+1}\Bigr)
 +(1-r)\sum_{j=1}^Lr^j\sum_{k=0}^{L-j}\sqrt{y_k}\,\pi_j(y_k),
\]
and the first term is $d$.

\emph{One grid.} Fix $j$ and $v=v_j$. Extend the inner sum to all
integers $k$; the terms with $k<0$ have $y_k<1\le v$ and add exactly
$d\sum_{i\ge1}2^{-i/2}=(1+\sqrt2)d$. Let $sv$, with $s\in[1,2)$, be the
first grid value $y_k$ at or above $v$. The grid values below it add
$\sqrt{sv}\sum_{i\ge1}2^{-i/2}=(1+\sqrt2)\sqrt{sv}$, and those from it on
add $\sqrt v\sum_{i\ge0}\sqrt{2^is}\,e^{-2v(\sqrt{2^is}-1)^2}$. So the
inner sum is at most $\sqrt{v_j}F_{v_j}(s)-(1+\sqrt2)d$.

\emph{The constant.} Let $C_1=89/20$, $C_2=17/4$ and $C_j=399/100$ for
$j\ge3$. Since $v_1=1+\log2>169/100$, $v_2>119/50$ and
$v_j\ge v_3>307/100$ for $j\ge3$, \cref{lem:upper-12-offset-cost} bounds
the inner sum by $\sqrt{v_j}C_j-(1+\sqrt2)d$. The subtracted terms add
$(1-r)(1+\sqrt2)d\sum_{j=1}^Lr^j=(1+r)(1-r^L)d$, which exceeds $d$ for
$L\ge3$ and pays for the first term. Hence
\[
 d\,\E[C_L\mid d,S]\le(1-r)\sum_{j\ge1}r^jC_j\sqrt{v_j}<5.19<\tfrac{26}{5}.
\]
The Lean proof bounds $\sqrt{v_j}$ by a rational number for $j\le8$ and by
$(v_j+9)/6$ for $j\ge9$, which gives about $5.185$; the sum itself is about
$5.182$. Here $\rho_N(d)=0$ and $\min\{\sqrt N,d^{-1}\}=d^{-1}$.
\end{proof}

\subsection{The second moment of one call}

Because $\Delta_\ell\in\{-1,0,1\}$,
\[
 \E[Y_L^2\mid d,S]=\sum_{\ell\le L}\frac{\Prob(\Delta_\ell\ne0)}{p_\ell},
 \qquad \frac1{p_\ell}=(2+\sqrt2)\,2^{\ell/2}.
\]
Bounding every probability by one gives
$(2+\sqrt2)\sum_{\ell\le L}2^{\ell/2}\le(2+\sqrt2)^2\sqrt N
=(6+4\sqrt2)\sqrt N<11.66\sqrt N$. This is
\eqref{eq:upper-12-moment} when $Nd^2\le1$, since
$\tfrac{13}{2}+\tfrac{129}{25}=11.66$. The rest of this section
proves $d\,\E[Y_L^2\mid d,S]\le13/2$, which is
\eqref{eq:upper-12-moment} when $Nd^2>1$.

For a dyadic cap $M$, write the two kinds of early stop of the
classifier with cap $M$ as
\begin{gather*}
 \mathcal E_M(d)=\sum_{n<M}e^{-2n(d+t_{n,M})^2},\qquad
 \mathcal D_M(d)=\sum_{n<M}\delta_{n,M}(d),\\
 \delta_{n,M}(d)=\begin{cases}e^{-2n(t_{n,M}-d)^2},&d\le t_{n,M},\\
 1,&d>t_{n,M},\end{cases}
\end{gather*}
summed over dyadic $n$, with $\mathcal E_{1/2}=0$; for $M=1$ both sums are
empty. By Hoeffding's inequality, for a coin with gap $d$ the sample mean
at the prefix $n$ crosses the threshold on the wrong side of $1/2$ with
probability at most $e^{-2n(d+t_{n,M})^2}$, and on the right side with
probability at most $\delta_{n,M}(d)$. The first kind is the
early-stopping part of the union bound \eqref{eq:error-sum}.

\begin{lemma}[Vote change]\label[lemma]{lem:upper-12-vote}
Let $m\ge1$ and let $X_1,\dots,X_{2m}$ be independent with a common
Bernoulli law. Then the majority votes after $m$ and after $2m$
observations differ with probability exactly $a(1-a)\le1/4$, where
$a=\Prob(\overline X_m>1/2)$:
\[
 \Prob\bigl(\ind\{\overline X_m>1/2\}\ne\ind\{\overline X_{2m}>1/2\}\bigr)
 =a(1-a)\le\tfrac14 .
\]
\end{lemma}

\begin{proof}
Let $U_1$ and $U_2$ count the successes among the first and the last $m$
observations. The votes are $\ind\{2U_1>m\}$ and $\ind\{U_1+U_2>m\}$, so
they differ exactly on
\[
 E=\{2U_1>m,\ U_1+U_2\le m\}\cup\{2U_1\le m,\ U_1+U_2>m\}.
\]
Let $E'$ be the same event with $U_1$ and $U_2$ exchanged. On the first
part of $E$, $U_2\le m-U_1<m/2$; on the second, $U_2>m-U_1\ge m/2$. So $E$
and $E'$ are disjoint, and $E\cup E'$ is the event that exactly one of
$2U_1>m$ and $2U_2>m$ holds. The counts are independent with the same law,
so $\Prob(E')=\Prob(E)$ and $2\Prob(E)=\Prob(E\cup E')=2a(1-a)$.
\end{proof}

\begin{lemma}\label[lemma]{lem:upper-12-level}
Let $\ell\ge1$ and $x=2^\ell d^2$.
\begin{enumerate}[label=(\roman*)]
\item $\Prob(\Delta_\ell\ne0)\le e^{-2x}+e^{-x}+\mathcal E_{2^\ell}(d)
 +\mathcal E_{2^{\ell-1}}(d)$.
\item $\Prob(\Delta_\ell\ne0)\le\tfrac14
 +\mathcal D_{2^{\ell-1}}(d)+\mathcal E_{2^{\ell-1}}(d)$.
\end{enumerate}
\end{lemma}

\begin{proof}
(i) The two classifiers disagree only if one of them is wrong. The union
bound \eqref{eq:error-sum} bounds the error of each by its final term
$e^{-2Md^2}$ plus $\mathcal E_M(d)$, with $Md^2=x$ and $x/2$.

(ii) Put $M=2^{\ell-1}\ge1$. The classifier with cap $M$ stops early with
probability at most $\mathcal D_M(d)+\mathcal E_M(d)$. Otherwise it
outputs $\ind\{\overline X_M>1/2\}$, and the classifier with cap $2M$,
whose thresholds are larger, has not stopped before $M$ either. If it
stops at $M$, it outputs the same side. If not, it outputs
$\ind\{\overline X_{2M}>1/2\}$. So disagreement without an early stop
requires the votes after $M$ and $2M$ observations to differ, which has
probability at most $1/4$ by \cref{lem:upper-12-vote} with $m=M$.
\end{proof}

Define the weight of level $\ell$, with $x_\ell=2^\ell d^2$, by
\[
 w_\ell=\begin{cases}
 e^{-2x_\ell}+e^{-x_\ell},& x_\ell\ge19/20,\\
 \tfrac14+\mathcal D_{2^{\ell-1}}(d),& x_\ell<19/20,\ \ell\ge1,\\
 1,&\ell=0.\end{cases}
\]
We call the value $19/20$ of $x_\ell$ the split. Since $x_1=2d^2\le1/2$,
the first case has $\ell\ge2$. By
\cref{lem:upper-12-level}, every level satisfies
$\Prob(\Delta_\ell\ne0)\le w_\ell+\mathcal E_{2^\ell}(d)
+\mathcal E_{2^{\ell-1}}(d)$. Since $d/p_\ell=(2+\sqrt2)\sqrt{x_\ell}$,
\begin{equation}\label{eq:upper-12-split}
 d\,\E[Y_L^2\mid d,S]\le(2+\sqrt2)\sum_{\ell\ge0}\sqrt{x_\ell}\,w_\ell
 +d\sum_{\ell\le L}\frac{\mathcal E_{2^\ell}(d)+\mathcal E_{2^{\ell-1}}(d)}
 {p_\ell}.
\end{equation}

\emph{Wrong-side early stops.} At the prefix $n=2^k$ of the cap
$M=2^\ell$, with $j=\ell-k$, $nt_{n,M}^2=v_j$, so
$n(d+t_{n,M})^2=(\sqrt{y_k}+\sqrt{v_j})^2\ge v_j+2\sqrt{v_j}\sqrt{y_k}$,
where $e^{-2v_j}=e^{-2}2^{-2j}$. The values $u_k=\sqrt{y_k}$ form a grid
of ratio $\sqrt2$, and for $\psi>0$,
\[
 \sum_ku_ke^{-\psi u_k}\le\frac{2+\sqrt2}{\psi},
\]
because $(1-r)z\,e^{-z}\le e^{-rz}-e^{-z}$ for $z\ge0$ and the right
side telescopes along the grid. With $d/p_\ell=(2+\sqrt2)2^{j/2}u_k$ and
$\psi=4\sqrt{v_j}\ge4\cdot\tfrac{13}{10}$, the terms of the larger classifier
add at most
$\tfrac{(2+\sqrt2)^2}{4\cdot13/10}e^{-2}\sum_{j\ge1}(\sqrt2/4)^j$. Those
of the smaller classifier add at most $\sqrt2$ times as much, because
$1/p_\ell=\sqrt2/p_{\ell-1}$. Since
$\sum_{j\ge1}(\sqrt2/4)^j=\sqrt2/(4-\sqrt2)$, the last sum in
\eqref{eq:upper-12-split} is at most
\[
 (1+\sqrt2)\,\frac{(2+\sqrt2)^2}{4\cdot13/10}\,e^{-2}\,
 \frac{\sqrt2}{4-\sqrt2}<0.401<\tfrac{41}{100}.
\]

\emph{Right-side early stops below the split.}
Let $M=2^{L'}$ with $L'\ge1$, and put $\xi=Md^2$. At the prefix $n=2^k$,
with $j=L'-k$, $nt_{n,M}^2=v_j$ and $2^kd^2=\xi2^{-j}$. If $\xi\le1/2$, then
$2^kd^2<1\le v_j$, so $d<t_{n,M}$ and
\[
 \mathcal D_M(d)\le\sum_{j\ge1}
 \exp\Bigl(-2\bigl(\sqrt{v_j}-\sqrt{\xi2^{-j}}\bigr)^2\Bigr),
\]
which increases with $\xi$.

The first value $x_\ell$ at or above the split lies in $[19/20,19/10)$.
Cut this interval into the four bands $[s_g,s_{g+1}]$, $g=0,1,2,3$, with
\[
 (s_0,s_1,s_2,s_3,s_4)=\bigl(\tfrac{19}{20},\tfrac{19}{16},\tfrac{57}{40},
 \tfrac{133}{80},\tfrac{19}{10}\bigr).
\]

\begin{lemma}\label[lemma]{lem:upper-12-drift}
Let $e_{g,i}$ be the entry of the following table for $i\le4$, and
$e_{g,i}=e_{g,4}$ for $i>4$. If $Md^2\le s_{g+1}2^{-i-2}$, then
$\mathcal D_M(d)\le e_{g,i}$.
\[
\begin{array}{c|ccccc}
 g & i=0 & i=1 & i=2 & i=3 & i=4\\\hline
 0 & 0.239064 & 0.156680 & 0.113601 & 0.089594 & 0.075413\\
 1 & 0.271967 & 0.173304 & 0.122528 & 0.094695 & 0.078491\\
 2 & 0.305008 & 0.189799 & 0.131249 & 0.099608 & 0.081419\\
 3 & 0.338195 & 0.206233 & 0.139823 & 0.104379 & 0.084231
\end{array}
\]
\end{lemma}

\begin{proof}
Since $s_{g+1}2^{-i-2}\le19/40<1/2$, it suffices to take
$\xi=s_{g+1}2^{-i-2}$ for $i\le4$. The Lean proofs bound the first six
terms with rational bounds on $\sqrt{v_j}$ and $\sqrt{\xi2^{-j}}$ and an
eight-term Taylor sum, and the terms with $j\ge7$ by $(31/100)^j$, which
add less than $1/2000$ for $\xi\le1/2$.
\end{proof}

\begin{lemma}\label[lemma]{lem:upper-12-offset-moment}
Let $\phi_0,\phi_1,\phi_2,\phi_3$ be $0.96$, $0.997$, $1.035$ and $1.073$.
For every $g$ and every $s\in[s_g,s_{g+1}]$,
\[
 \phi_g\sqrt s+\sum_{i\ge0}\sqrt{2^is}\,\bigl(e^{-2^{i+1}s}+e^{-2^is}\bigr)
 \le\tfrac{89}{50}.
\]
\end{lemma}

\begin{proof}
The Lean proof splits the four bands into $17$ pieces in all, chosen and
checked in exact rational arithmetic. On each piece it bounds the first six
terms of the series by the monotonicity of each factor and an eight-term
Taylor lower bound for the exponential, and the remaining terms by
$1/1000$. The maximum of the left side is about $1.762$.
\end{proof}

\begin{proof}[Proof of \eqref{eq:upper-12-moment}]
It remains to bound the first sum in \eqref{eq:upper-12-split}. Let
$\ell_*$ be the first level with $x_{\ell_*}\ge19/20$ and
$s=x_{\ell_*}\in[19/20,19/10)$, in the band $g$. Then $\ell_*\ge2$.

\emph{Levels below $\ell_*$.} Write $\ell=\ell_*-1-i$, so that
$\sqrt{x_\ell}=r^{i+1}\sqrt s$. For $\ell\ge1$,
$x_{\ell-1}=s2^{-i-2}\le s_{g+1}2^{-i-2}$, so
$w_\ell\le\tfrac14+e_{g,i}$ by \cref{lem:upper-12-drift}. Level $0$,
with $i=\ell_*-1$, has weight one. So the levels below $\ell_*$ add at
most
\[
 \sqrt s\Bigl(\sum_{i<\ell_*-1}r^{i+1}\bigl(\tfrac14+e_{g,i}\bigr)+r^{\ell_*}\Bigr)
 \le\phi_g\sqrt s,
\]
which the Lean proof checks for every value of $\ell_*$ from the table of
\cref{lem:upper-12-drift}.

\emph{Levels from $\ell_*$ on.} The level $\ell_*+i$ adds
$\sqrt{2^is}\,(e^{-2^{i+1}s}+e^{-2^is})$.

By \cref{lem:upper-12-offset-moment} the first sum in
\eqref{eq:upper-12-split} is at most $89/50$. Hence
\[
 d\,\E[Y_L^2\mid d,S]\le(2+\sqrt2)\tfrac{89}{50}+\tfrac{41}{100}
 <6.49<\tfrac{13}{2}.\qedhere
\]
\end{proof}

\subsection{What Lean checks and what remains}\label{sec:upper-12-scope}

The Lean proof of \cref{thm:upper-12} builds on \texttt{BlindUpper},
which supplies the cutoff estimator, the calibrated batch and its budget,
on \texttt{UpperWork}, which supplies \eqref{eq:upper-12-survival}, on
\texttt{UpperMoment}, which supplies the weights $1/p_\ell$ and their sums,
and on \texttt{UpperMixture}, which supplies \eqref{eq:scaling} in the form
used here. Six groups of modules, which also hold the audited stages $19$,
$17$, $15$, $13$ and $12$ of \cref{sec:verification}, add the rest:
\begin{itemize}
\item \texttt{UpperCost} proves the first two bounds of
 \cref{lem:upper-12-offset-cost} and the lemmas of the reordering by $j$.
\item \texttt{UpperDisagreement} proves part (i) of
 \cref{lem:upper-12-level} and the bound $41/100$ on the wrong-side early
 stops, and defines the regime term $\rho_N$.
\item \texttt{UpperSign} proves the drift series bounds for $\xi\le3/8$,
 which \texttt{UpperQuarter} extends to $\xi\le1/2$.
\item \texttt{UpperQuarter} proves the bound $1/4$ of
 \cref{lem:upper-12-vote}, whose proof establishes the identity, part (ii) of
 \cref{lem:upper-12-level},
 \cref{lem:upper-12-drift,lem:upper-12-offset-moment} and
 \eqref{eq:upper-12-moment}, with its regime term.
\item \texttt{UpperCallCost} proves the third bound of
 \cref{lem:upper-12-offset-cost}, the constant $26/5$, the reordering by
 $j$ with these offsets, and \eqref{eq:upper-12-cost}, with its regime
 term.
\item \texttt{UpperScale} proves \cref{lem:upper-12-batch} and
 \cref{thm:upper-12}. The batch machine of \texttt{BlindUpper} marks a
 call when its uniform seed mark $u$ satisfies $u<C_L/2^L$; run on seeds
 whose marks are multiplied by $\lambda$, the same machine marks it when
 $u<C_L/(\lambda2^L)$, and so realizes the batch at scale $\lambda2^L$ with
 the same calls.
\end{itemize}
The two targets were stated after the proof and are pinned by hash. The
audit \path{tools/verify_upper_scale.py} restates them and twelve components
with their types written out and reports the axioms of all 90 public
declarations of the last two groups (record:
\path{results/upper-scale-20261119/}). The confidence bound of
\cref{thm:main} and \cref{app:hard-budget,app:sublinear-power,app:observed-upper}
keep the family of \cref{sec:minimax-proof} and their stated constants.

Lean does not check the approximate values in the text ($5.182$, $5.185$,
$1.762$, $12.5$, and the bounds $5.19$, $0.401$, $6.49$, for which it checks
$26/5$, $41/100$, $13/2$) or the remarks after the proof of \cref{thm:upper-12}.
A numerical study in \path{research/2026-11-19-upper-constant-12/} puts the
worst ratio of the family of \cref{thm:upper-12} near $5.4$, and the exact
constants of one call near $2.2$ for the second moment and $3.6$ for the cost,
against $13/2$ and $26/5$ here. Below the split the bound $\tfrac14+e_{g,i}$
is about $0.3$ to $0.6$, and the cost bound takes the survival probability to
be one below each threshold.
